% ===========================================================================
%  Theorem system - paper implementation (look "paper")
% ---------------------------------------------------------------------------
%  No boxes, no colour: plain amsthm environments as used by every physics and
%  mathematics journal.
%
%    plain style       Theorem, Lemma, Corollary, Proposition, Fact, Claim
%                      bold run-in head, italic statement
%    definition style  Definition, Example, Exercise, Homework
%                      bold run-in head, upright text
%    remark style      Remark, unnumbered Example, unnumbered Fact
%                      italic head, upright text
%    proof             amsthm proof environment with the QED square
%
%  Numbering matches the boxed look: one shared counter, numbered within the
%  section, so switching looks never renumbers anything.
%
%  Selected by style/theorems.tex, which documents the commands. Every command
%  defined in style/theorems-boxed.tex must also be defined here.
% ===========================================================================

% ---------------------------------------------------------------------------
%  Environments
% ---------------------------------------------------------------------------
\theoremstyle{plain}
\newtheorem{papertheorem}{Theorem}[section]
\newtheorem{paperlemma}[papertheorem]{Lemma}
\newtheorem{papercorollary}[papertheorem]{Corollary}
\newtheorem{paperproposition}[papertheorem]{Proposition}
\newtheorem{paperfact}[papertheorem]{Fact}
\newtheorem*{paperclaim}{Claim}

\theoremstyle{definition}
\newtheorem{paperdefinition}[papertheorem]{Definition}

\theoremstyle{remark}
\newtheorem*{paperexample}{Example}
\newtheorem*{paperremark}{Remark}
\newtheorem*{paperfactstar}{Fact}

% Compatibility aliases for the environment names of the boxed look, so that
% content written against either implementation compiles under both.
\newenvironment{mydefinition}[2]{\begin{paperdefinition}[#1]}{\end{paperdefinition}}
\newenvironment{mytheorem}[2]{\begin{papertheorem}[#1]}{\end{papertheorem}}
\newenvironment{mylemma}[2]{\begin{paperlemma}[#1]}{\end{paperlemma}}
\newenvironment{mycorollary}[2]{\begin{papercorollary}}{\end{papercorollary}}
\newenvironment{myproposition}[2]{\begin{paperproposition}}{\end{paperproposition}}
\newenvironment{example}{\begin{paperexample}}{\end{paperexample}}
\newenvironment{remark}{\begin{paperremark}}{\end{paperremark}}
\newenvironment{factenv}{\begin{paperfactstar}}{\end{paperfactstar}}

% Proof blocks. The boxed look gives each box type its own coloured rule; on
% paper they are all the one standard proof environment.
\newenvironment{projectproof}[1][\proofname]{\begin{proof}[#1]}{\end{proof}}
\newenvironment{thmpf}{\begin{proof}}{\end{proof}}
\newenvironment{lempf}{\begin{proof}}{\end{proof}}
\newenvironment{corpf}{\begin{proof}}{\end{proof}}
\newenvironment{proppf}{\begin{proof}}{\end{proof}}
\newenvironment{clmpf}{\begin{proof}}{\end{proof}}

% ---------------------------------------------------------------------------
%  Definition - the title is still indexed automatically
% ---------------------------------------------------------------------------
\NewDocumentCommand{\defn}{o m+m}{%
    \IfNoValueTF{#1}{\index{#2}}{%
        \def\defn@key{#1}\def\defn@nosig{-}%
        \ifx\defn@key\defn@nosig\else\index{#1}\fi%
    }%
    \IfBlankTF{#2}
        {\begin{paperdefinition}#3\end{paperdefinition}}
        {\begin{paperdefinition}[#2]#3\end{paperdefinition}}%
}

\NewDocumentCommand{\defnr}{o mm+m}{%
    \IfNoValueTF{#1}{\index{#2}}{%
        \def\defn@key{#1}\def\defn@nosig{-}%
        \ifx\defn@key\defn@nosig\else\index{#1}\fi%
    }%
    \IfBlankTF{#2}
        {\begin{paperdefinition}\label{defn:#3}#4\end{paperdefinition}}
        {\begin{paperdefinition}[#2]\label{defn:#3}#4\end{paperdefinition}}%
}

% ---------------------------------------------------------------------------
%  Theorem
% ---------------------------------------------------------------------------
\NewDocumentCommand{\thm}{m+m}{%
    \IfBlankTF{#1}
        {\begin{papertheorem}#2\end{papertheorem}}
        {\begin{papertheorem}[#1]#2\end{papertheorem}}%
}

\NewDocumentCommand{\thmr}{mm+m}{%
    \IfBlankTF{#1}
        {\begin{papertheorem}\label{thm:#2}#3\end{papertheorem}}
        {\begin{papertheorem}[#1]\label{thm:#2}#3\end{papertheorem}}%
}

\NewDocumentCommand{\thmp}{m+m+m}{%
    \thm{#1}{#2}
    \begin{proof}#3\end{proof}%
}

\NewDocumentCommand{\thmpr}{mm+m+m}{%
    \thmr{#1}{#2}{#3}
    \begin{proof}#4\end{proof}%
}

% ---------------------------------------------------------------------------
%  Lemma
% ---------------------------------------------------------------------------
\NewDocumentCommand{\lem}{m+m}{%
    \IfBlankTF{#1}
        {\begin{paperlemma}#2\end{paperlemma}}
        {\begin{paperlemma}[#1]#2\end{paperlemma}}%
}

\NewDocumentCommand{\lemr}{mm+m}{%
    \IfBlankTF{#1}
        {\begin{paperlemma}\label{lem:#2}#3\end{paperlemma}}
        {\begin{paperlemma}[#1]\label{lem:#2}#3\end{paperlemma}}%
}

\NewDocumentCommand{\lemp}{m+m+m}{%
    \lem{#1}{#2}
    \begin{proof}#3\end{proof}%
}

\NewDocumentCommand{\lempr}{mm+m+m}{%
    \lemr{#1}{#2}{#3}
    \begin{proof}#4\end{proof}%
}

% ---------------------------------------------------------------------------
%  Corollary (no title)
% ---------------------------------------------------------------------------
\NewDocumentCommand{\cor}{+m}{%
    \begin{papercorollary}#1\end{papercorollary}%
}

\NewDocumentCommand{\corr}{m+m}{%
    \begin{papercorollary}\label{cor:#1}#2\end{papercorollary}%
}

\NewDocumentCommand{\corp}{+m+m}{%
    \begin{papercorollary}#1\end{papercorollary}
    \begin{proof}#2\end{proof}%
}

% ---------------------------------------------------------------------------
%  Proposition (no title)
% ---------------------------------------------------------------------------
\NewDocumentCommand{\prop}{+m}{%
    \begin{paperproposition}#1\end{paperproposition}%
}

\NewDocumentCommand{\propr}{m+m}{%
    \begin{paperproposition}\label{prop:#1}#2\end{paperproposition}%
}

\NewDocumentCommand{\propp}{+m+m}{%
    \begin{paperproposition}#1\end{paperproposition}
    \begin{proof}#2\end{proof}%
}

% ---------------------------------------------------------------------------
%  Claim (unnumbered)
% ---------------------------------------------------------------------------
\NewDocumentCommand{\clm}{m+m}{%
    \IfBlankTF{#1}
        {\begin{paperclaim}#2\end{paperclaim}}
        {\begin{paperclaim}[#1]#2\end{paperclaim}}%
}

\NewDocumentCommand{\clmp}{m+m+m}{%
    \clm{#1}{#2}
    \begin{proof}#3\end{proof}%
}

% ---------------------------------------------------------------------------
%  Fact
% ---------------------------------------------------------------------------
\NewDocumentCommand{\fact}{+m}{%
    \begin{paperfact}#1\end{paperfact}%
}

\NewDocumentCommand{\factb}{+m}{%
    \begin{paperfactstar}#1\end{paperfactstar}%
}

% ---------------------------------------------------------------------------
%  Proof, example, remark
% ---------------------------------------------------------------------------
\NewDocumentCommand{\pf}{+m}{%
    \begin{proof}#1\end{proof}%
}

\NewDocumentCommand{\ex}{+m}{%
    \begin{paperexample}#1\end{paperexample}%
}

\NewDocumentCommand{\rmkb}{+m}{%
    \begin{paperremark}#1\end{paperremark}%
}

% Inline remark: an italic aside inside running text
\NewDocumentCommand{\rmk}{+m}{%
    {\itshape\color{projectrmk}#1}%
}
